\chapter*{Conclusion Générale}
\addcontentsline{toc}{chapter}{Conclusion Générale}
\markboth{Conclusion Générale}{}

La réalisation de ce projet, mené au sein de la Fondation Djagora en collaboration avec l'Institut Supérieur de Gestion Industrielle de Sfax (ISGIS), a permis de concevoir et de déployer LogiLink, une plateforme intelligente d'adéquation entre la formation et l'emploi dédiée au secteur du transport et de la logistique. Ce système apporte une réponse concrète au défi de l'alignement des compétences académiques avec les exigences du marché du travail, tout en accélérant l'insertion professionnelle des futurs diplômés.

\noindent Notre contribution au sein de ce projet s'est concentrée sur le développement de l'intelligence analytique de la plateforme à travers l'implémentation de quatre moteurs importants : la génération automatisée de CV, le calcul du score ATS, le calcul du score d'employabilité ainsi que le calcul du score d'employabilité prévisionnelle. L'intégration de ces modules offre une valeur opérationnelle précise : elle permet à l'étudiant de suivre ses scores et de générer son CV, offre aux sociétés la visibilité nécessaire sur le score d'employabilité des candidats, et fournit à l'administration une vue d'ensemble complète pour assurer un pilotage académique et professionnel.

\noindent Sur le plan de la gestion de projet, l'adoption de la méthodologie Agile Scrum a garanti une orchestration rigoureuse des cycles de développement, assurant une communication fluide entre les parties prenantes et le maintien d'une dynamique de livraison continue. L'architecture technique, hautement modulable, repose sur les frameworks Angular et NestJS, associés à des microservices Python dédiés aux algorithmes de scoring et d'intelligence artificielle.

\noindent En termes de perspectives, l'évolution de ce travail réside dans l'automatisation de la mise à jour des référentiels métiers et des calculs de scores en se basant directement sur les offres d'emploi réelles et les critères définis par les sociétés partenaires sur la plateforme. À plus long terme, l'évolution du système passera par l'amélioration continue des performances globales de la plateforme, notamment à travers l'optimisation du module de scoring et le renforcement de la pertinence des résultats d'évaluation.

\noindent En conclusion, ce projet représente bien plus qu'une simple réalisation technique. Il marque une étape clé dans notre parcours en consolidant notre expertise technique en Business Intelligence, en développement full-stack et en intelligence artificielle. Au-delà de son aspect technologique, la solution développée apporte une contribution concrète au monde académique en permettant d'optimiser l'employabilité des étudiants tout en renforçant l'adéquation entre la formation universitaire et le marché de l'emploi dans le secteur du Transport et de la Logistique. 
\newpage
